% ============================================================
%  国赛（CUMCM）论文骨架 — cumcmthesis v2.9
%  2026-09-09 入库（独立审计修复 P0-1/P1-3：此前模板分支声称 main.tex 存在但实际缺失）
%  编译链：cd paper/ && xelatex main.tex && xelatex main.tex（thebibliography 内联文献，无需 bibtex）
%  ⛔ 表格规范：三线表用 \toprule/\midrule/\bottomrule（cls 已修复加载顺序，booktabs 最后加载）
%  ⛔ 勿在 main.tex 重复加载 subcaption/float/graphicx/booktabs（cls 已含）
%  ⛔ 勿加 \usepackage{geometry}（页边距由 cls 统一设定）
% ============================================================
\documentclass[withoutpreface,bwprint]{cumcmthesis}

\begin{document}

\title{论文标题（20 字以内，准确体现研究对象与模型方法）}
\date{\today}

% ---------- 摘要：最后写！必须引用各章真实数值 ----------
\begin{abstract}
针对……问题，建立了……模型（第一段：总述方法与结论）。
关于问题一，基于……方法，求得……（给出关键数值结果与检验结论）。
关于问题二，……
关于问题三，……
本文的创新点/局限性……
\keywords{关键词一\quad 关键词二\quad 关键词三\quad 关键词四\quad 关键词五}
\end{abstract}

\tableofcontents

\section{问题重述}
% 1.1 问题背景（简述，勿抄题面）；1.2 需要解决的问题（逐条列出）
\subsection{问题背景}
……
\subsection{需要解决的问题}
……

\section{问题分析}
% 每个子问题：数据特征 → 方法选择 → 可行性论证
针对该问题的求解范式，文献\cite{ref1}给出了经典建模思路；针对大规模情形，
启发式算法与精确算法的取舍见文献\cite{ref2}。

\section{模型假设与符号说明}
\subsection{模型假设}
\begin{enumerate}
  \item 假设一（每条假设给出依据）。
  \item 假设二……
\end{enumerate}
\subsection{符号说明}
% 三线表示例（booktabs；长表用 longtable + \endfirsthead/\endhead）
\begin{table}[htbp]
\centering
\caption{主要符号说明}
\begin{tabular}{clc}
\toprule
符号 & 含义 & 单位 \\
\midrule
$x_i$ & 第 $i$ 个决策变量 & --- \\
$W$   & 资源总量上限 & 吨 \\
\bottomrule
\end{tabular}
\end{table}

\section{模型的建立与求解}
% 每个子问题一小节：模型（公式编号）→ 求解算法 → 结果表/图 → 结果分析
\subsection{问题一模型的建立与求解}
……

\section{结果的分析与检验}
% 灵敏度分析 / 稳健性检验 / 误差分析——每个模型至少一张灵敏度分析表
……

\section{模型的评价与推广}
% 优点、缺点、改进方向、推广应用——具体且有量词，禁空话
……

% ---------- 参考文献：inline thebibliography（国赛规范） ----------
\begin{thebibliography}{99}
\bibitem{ref1} 作者. 书名[M]. 出版地: 出版社, 年份.
\bibitem{ref2} 作者. 论文题目[J]. 期刊名, 卷(期): 起止页码, 年份.
\end{thebibliography}

\appendix
\section{源程序清单}
% 只放核心代码；用 listing 或 \input 引入
\section{支撑文件列表}
% 按 comp-code 交付清单列出全部代码与产出文件

\end{document}
